\chapter{The Agent Layer}
\label{ch:agents}

This is where the design lives. Four classes: one abstract base, three specialists that differ only in a prompt and a filter, and an orchestrator that has all the control flow.

\section{BaseAgent: the template}

\code{BaseAgent} is 63 lines and defines the entire agent contract through two abstract methods and one concrete pipeline.

\begin{lstlisting}[style=py,firstnumber=5]
class BaseAgent(ABC):
    def __init__(self, name: str, rag_retriever, llm_client,
                 temperature: float = 0.3, max_tokens: int = 3000):
        self.name          = name
        self.rag_retriever = rag_retriever
        self.llm_client    = llm_client
        self.temperature   = temperature
        self.max_tokens    = max_tokens

    @abstractmethod
    def get_system_prompt(self) -> str: ...

    @abstractmethod
    def get_category_filter(self) -> Optional[str]: ...
\end{lstlisting}

A subclass supplies a personality (the system prompt) and a slice of the corpus (the category filter). Everything else is inherited. Adding a fourth specialist is therefore a 15-line file --- see Chapter~\ref{ch:extend}.

\subsection{process(): the three-step pipeline}

\begin{lstlisting}[style=py,firstnumber=59]
    def process(self, query: str, additional_info: Dict = None,
                conversation_history: list = None) -> Dict:
        source_filter = additional_info.get('source_filter') if additional_info else None
        context, sources = self.retrieve_context_and_sources(query, source_filter=source_filter)
        response = self.generate_response(query, context, additional_info, conversation_history)
        return {'agent': self.name, 'response': response,
                'sources': sources, 'query': query}
\end{lstlisting}

Retrieve, generate, return both the prose and the chunks it came from. Returning \code{sources} alongside \code{response} is what lets the orchestrator build a citation list without re-running retrieval.

\subsection{Retrieval with the agent's own filter}

\begin{lstlisting}[style=py,firstnumber=21]
    def retrieve_context_and_sources(self, query: str, n_results: int = 5,
                                     source_filter: Optional[str] = None):
        category_filter = self.get_category_filter()
        chunks = self.rag_retriever.retrieve(query=query, n_results=n_results,
                                             category_filter=category_filter,
                                             source_filter=source_filter)
        context_parts = []
        for i, chunk in enumerate(chunks, 1):
            metadata = chunk.get('metadata', {})
            source = (f"[Source: {metadata.get('section', 'Unknown')} - "
                      f"{metadata.get('chapter', 'Unknown')}]")
            context_parts.append(f"--- Context {i} {source} ---")
            context_parts.append(chunk.get('text', ''))
            context_parts.append("")
        return "\n".join(context_parts), chunks
\end{lstlisting}

The category filter comes from the subclass and is never overridden by the caller. That is the mechanism behind the whole multi-agent premise: three agents asking the same corpus the same question get three different answers because each one is only allowed to see a third of it.

\subsection{Passing structured findings between agents}

\code{generate\_response} folds an \code{additional\_info} dictionary into the query as labelled lines, after removing the one key that is meant for retrieval rather than for the model:

\begin{lstlisting}[style=py,firstnumber=43]
        # Ensure 'source_filter' is not included in the LLM's additional_info
        llm_additional_info = additional_info.copy() if additional_info else {}
        llm_additional_info.pop('source_filter', None)

        if llm_additional_info:
            info_str = "\n".join([f"{key}: {value}"
                                  for key, value in llm_additional_info.items()])
            enhanced_query = f"{query}\n\nAdditional Information:\n{info_str}"
        else:
            enhanced_query = query
\end{lstlisting}

So the treatment agent's actual query text ends up looking like:

\begin{lstlisting}[style=plain]
I have acid reflux and skin rashes. What should I do?

Additional Information:
Prakriti: **Provisional Constitutional Assessment**: Pitta-dominant, high confidence.
  The user's ...
Dosha Imbalance: **Provisional Diagnosis (Vikriti Nirnaya)**: Pitta Prakopa with
  Tikshnagni. Amlapitta is indicated by ...
\end{lstlisting}

Prose is passed between agents, not structured data. There is no schema, no parsing, no validation --- the downstream agent reads its predecessor's report the way a junior clinician reads a colleague's notes. That keeps the code short and makes the handoff lossy in a way that is hard to measure.

\begin{note}[title={One retrieval too many}]
\code{process()} calls \code{retrieve\_context\_and\_sources()} and then passes the resulting context into \code{generate\_response()}. But \code{generate\_response()} also retrieves when \code{context is None}, so it carries a second, unreachable retrieval path. The dead branch is harmless; it just means two functions know how to fetch context.
\end{note}

\section{The three specialists}

\begin{table}[H]
\centering
\small
\begin{tabularx}{\textwidth}{@{}l l c l X@{}}
\toprule
\textbf{Class} & \textbf{\code{name}} & \textbf{Temp} & \textbf{Filter} & \textbf{Persona in the system prompt} \\
\midrule
\code{PrakritiAgent} & Prakriti Assessor & 0.3 & \code{prakriti} & Senior professor of Sharira Vijnana \\
\code{DoshaAgent} & Dosha Imbalance Detector & 0.3 & \code{vikriti} & Senior clinician in Vikriti Vijnana \\
\code{TreatmentAgent} & Treatment Recommender & 0.4 & \code{treatment} & Senior physician designing Chikitsa Krama \\
\bottomrule
\end{tabularx}
\caption{The three specialists. Diagnosis runs cooler than therapy design, which is given slightly more latitude at 0.4.}
\end{table}

Each class is under 40 lines and most of that is the prompt string. The whole of \fpath{src/agents/prakriti\_agent.py} outside the prompt:

\begin{lstlisting}[style=py]
class PrakritiAgent(BaseAgent):
    def __init__(self, rag_retriever, llm_client):
        super().__init__(name="Prakriti Assessor", rag_retriever=rag_retriever,
                         llm_client=llm_client, temperature=0.3)

    def get_system_prompt(self) -> str:
        return """You are a senior Ayurvedic professor specializing in ..."""

    def get_category_filter(self) -> Optional[str]:
        return "prakriti"
\end{lstlisting}

\subsection{What each prompt demands}

The prompts are the specification, so they are worth reading closely. All three share four instructions --- scholarly tone, classical citation, Sanskrit terminology with parenthetical glosses, and a fixed output structure --- and then diverge on what they must produce.

\textbf{Prakriti Assessor} must return a constitutional call with a confidence level, a per-dosha breakdown naming the specific Gunas, and a trait-to-Guna correlation table. The prompt supplies the table format inline:

\begin{lstlisting}[style=plain]
| User Trait  | Assessed Dosha | Corresponding Guna |
|-------------|----------------|--------------------|
| Dry Skin    | Vata           | Ruksha (Dry)       |
| Irritability| Pitta          | Tikshna (Sharp)    |
\end{lstlisting}

\textbf{Dosha Imbalance Detector} must go past naming a dosha and describe the pathogenesis: the nature of the vitiation (\emph{Vata Prakopa}, \emph{Pitta Kshaya}), which tissues and channels are involved (\emph{Dhatus}, \emph{Srotas}), and the state of the digestive fire (\emph{Mandagni}, \emph{Tikshnagni}).

\textbf{Treatment Recommender} has the longest prompt and the only safety-critical clause. It is instructed to lead with the treatment principle, then split recommendations into purification and palliative care, and it carries a mandatory string:

\begin{lstlisting}[style=plain]
Advanced procedures like Vamana, Virechana, and Basti *must* include the
statement: "To be performed under the direct supervision of a qualified Vaidya."
\end{lstlisting}

It also branches on which Sthana the user asked about, which is the other half of the source-filtering feature:

\begin{lstlisting}[style=plain]
* Kalpa Sthana:    Focus on preparation and properties of specific herbs
                   for Panchakarma (e.g., Madanaphala, Vacha).
* Siddhi Sthana:   Focus on management of Panchakarma procedures and
                   post-therapy care (Paschat Karma), like Samsarjana Krama.
* Chikitsa Sthana: Focus on disease-specific formulations and protocols.
\end{lstlisting}

\begin{warn}[title={Safety enforcement is prompt-level only}]
The supervision warning for Vamana, Virechana and Basti is an instruction to the model, not a post-processing check. Nothing in the code verifies that the string appears in the output before it reaches the user. The synthesis prompt asks the orchestrator to re-check it, which makes it twice as likely and still not guaranteed. A regex assertion over the final text before display would close this.
\end{warn}

\section{The orchestrator}

\subsection{Routing}

\code{analyze\_query} is three \code{any()} calls over lowercased substrings:

\begin{lstlisting}[style=py,firstnumber=27]
    def analyze_query(self, query: str) -> Dict:
        query_lower = query.lower()
        needs_prakriti = any(term in query_lower for term in
            ['constitution', 'prakriti', 'vata', 'pitta', 'kapha', 'body type'])
        needs_dosha = any(term in query_lower for term in
            ['symptom', 'problem', 'pain', 'disorder', 'disease', 'sick', 'imbalance'])
        needs_treatment = any(term in query_lower for term in
            ['treatment', 'remedy', 'cure', 'help', 'diet', 'food', 'herb', 'medicine'])

        # This no longer defaults to activating all agents. If no keywords match,
        # it correctly signals a general, non-Ayurvedic query.
        return {'prakriti': needs_prakriti, 'dosha': needs_dosha,
                'treatment': needs_treatment}
\end{lstlisting}

\begin{table}[H]
\centering
\small
\begin{tabularx}{\textwidth}{@{}l X@{}}
\toprule
\textbf{Agent} & \textbf{Trigger words} \\
\midrule
Prakriti & constitution, prakriti, vata, pitta, kapha, body type \\
Dosha & symptom, problem, pain, disorder, disease, sick, imbalance \\
Treatment & treatment, remedy, cure, help, diet, food, herb, medicine \\
\bottomrule
\end{tabularx}
\caption{The complete routing vocabulary. Twenty words decide the entire execution path.}
\end{table}

The comment in the code marks a real design change --- an earlier version fired all three agents when nothing matched, and the git history shows the fix landing as \emph{``feat: Implement smart query routing in OrchestratorAgent''}. The current behaviour costs one model call for a general question instead of four.

Two things follow from a substring match. First, \code{'help'} is a treatment keyword, so ``help'' anywhere in a sentence activates the treatment agent --- ``can you help me understand Ayurveda?'' takes the slow path and gets a treatment protocol it did not ask for. Second, a genuine clinical description that avoids all twenty words takes the fast path with no retrieval at all: ``I wake at 3am with a racing mind and my joints crack'' contains none of them.

\subsection{The two paths}

\begin{lstlisting}[style=py,firstnumber=149]
    def simple_query(self, query: str, conversation_history: List[Dict] = None) -> Dict:
        agent_activation   = self.analyze_query(query)
        is_ayurvedic_query = any(agent_activation.values())

        if is_ayurvedic_query:
            # SLOW PATH: full multi-agent process
            return self.process_query(query, agent_activation, conversation_history)
        else:
            # FAST PATH: bypass agents for a direct answer
            system_prompt = ("You are AyurMind, a friendly and knowledgeable "
                             "Ayurvedic assistant. Answer general questions "
                             "conversationally and helpfully. ...")
            response = self.llm_client.generate(
                prompt=query, system_prompt=system_prompt,
                temperature=0.4, max_tokens=2000,
                conversation_history=conversation_history)
            return {'query': query, 'final_response': response,
                    'retrieved_sources': [], 'agent_activation': agent_activation}
\end{lstlisting}

Note the tone switch. The fast path is explicitly ``friendly''; the slow path's synthesis prompt explicitly forbids friendliness. The same system speaks in two registers depending on which words appeared in the question.

\subsection{The cascade}

\code{process\_query} runs the activated agents in a fixed order and feeds each one's output forward.

\begin{figure}[H]
\centering
\begin{tikzpicture}[node distance=9mm]
  \node[blk=deepgreen, minimum width=28mm] (p) {Prakriti Agent};
  \node[blk=deepgreen, right=26mm of p, minimum width=28mm] (d) {Dosha Agent};
  \node[blk=deepgreen, right=26mm of d, minimum width=28mm] (t) {Treatment Agent};
  \node[blk=clay, below=16mm of d, minimum width=44mm] (s) {Synthesis (orchestrator)};

  \draw[arb=deepgreen] (p) -- node[lbl,above,align=center]{\code{Prakriti}\\\code{Assessment}} (d);
  \draw[arb=deepgreen] (d) -- node[lbl,above,align=center]{\code{Dosha}\\\code{Imbalance}} (t);
  \draw[arb=deepgreen] (p.south) to[out=-70,in=170] node[lbl,left,yshift=-3mm]{\code{Prakriti}} (s.west);
  \draw[arb=clay] (d) -- (s);
  \draw[arb=deepgreen] (t.south) to[out=-110,in=10] (s.east);

  \node[db=plum, above=11mm of p, minimum width=22mm] (r1) {\scriptsize prakriti};
  \node[db=plum, above=11mm of d, minimum width=22mm] (r2) {\scriptsize vikriti};
  \node[db=plum, above=11mm of t, minimum width=22mm] (r3) {\scriptsize treatment};
  \draw[arb=plum] (r1) -- (p);
  \draw[arb=plum] (r2) -- (d);
  \draw[arb=plum] (r3) -- (t);
\end{tikzpicture}
\caption{Information cascade. Each agent retrieves from its own category slice, and each downstream agent additionally receives the prose written by the ones before it.}
\end{figure}

\begin{lstlisting}[style=py,firstnumber=52]
        if agent_activation['prakriti']:
            prakriti_result = self.prakriti_agent.process(
                query, base_additional_info.copy(),
                conversation_history=conversation_history)
            results['prakriti'] = prakriti_result['response']
            if 'sources' in prakriti_result:
                retrieved_sources.extend(prakriti_result['sources'])

        if agent_activation['dosha']:
            dosha_additional_info = base_additional_info.copy()
            if 'prakriti' in results:
                dosha_additional_info['Prakriti Assessment'] = results['prakriti']
            dosha_result = self.dosha_agent.process(
                query, dosha_additional_info, conversation_history)
            results['dosha'] = dosha_result['response']
            ...

        if agent_activation['treatment']:
            treatment_additional_info = base_additional_info.copy()
            if 'prakriti' in results:
                treatment_additional_info['Prakriti'] = results['prakriti']
            if 'dosha' in results:
                treatment_additional_info['Dosha Imbalance'] = results['dosha']
            ...
\end{lstlisting}

The ordering is not arbitrary --- it mirrors clinical sequence. You establish the baseline constitution, then read the deviation from it, then prescribe against the deviation. The guards (\code{if 'prakriti' in results}) mean a partially activated run still works: a treatment-only query gets a treatment plan with no constitutional context, which is exactly what a bare question like ``what foods help sleep?'' deserves.

Sources from all activated agents are pooled and deduplicated by chunk ID:

\begin{lstlisting}[style=py,firstnumber=79]
        # Remove duplicate sources based on 'id'
        unique_sources = list({s['id']: s for s in retrieved_sources}.values())
\end{lstlisting}

A dictionary comprehension keyed on ID keeps the last occurrence of each. With all three agents active the pool is up to 15 chunks before dedup; the category filters make overlap unlikely, so it usually stays near 15.

\subsection{Source extraction}

If the user names a classical section, retrieval is narrowed to it:

\begin{lstlisting}[style=py,firstnumber=14]
    def _extract_requested_sources(self, query: str) -> Optional[str]:
        """Extracts requested sources like 'Siddhi Sthana' using regex."""
        match = re.search(r"(?:from|in)\s+((?:\w+\s+)?(?:Sthana|Samhita|Kalpa))",
                          query, re.IGNORECASE)
        if match:
            return match.group(1)
        # A simpler check if the above fails
        if "siddhi sthana" in query.lower():
            return "Siddhi Sthana"
        if "kalpa sthana" in query.lower():
            return "Kalpa Sthana"
        return None
\end{lstlisting}

The regex requires a preceding \code{from} or \code{in}, so ``according to the Siddhi Sthana'' misses the pattern and is caught by the hard-coded fallback beneath it --- which covers only two of the eight sections. Whatever it extracts goes into \code{additional\_info['source\_filter']} and flows down to every activated agent.

As Chapter~5 explains, the extracted string never matches the stored \code{section} value, so this feature currently narrows retrieval to nothing rather than to the requested book. It is a two-line fix in the metadata schema and worth doing, because the treatment prompt already knows how to behave differently per Sthana.

\subsection{Synthesis}

The last step assembles the specialists' reports plus a source list and asks for one document.

\begin{lstlisting}[style=py,firstnumber=87]
        synthesis_context = "Agent Analyses:\n\n"
        if 'prakriti' in agent_results:
            synthesis_context += f"CONSTITUTIONAL ASSESSMENT:\n{agent_results['prakriti']}\n\n"
        if 'dosha' in agent_results:
            synthesis_context += f"IMBALANCE ANALYSIS:\n{agent_results['dosha']}\n\n"
        if 'treatment' in agent_results:
            synthesis_context += f"TREATMENT RECOMMENDATIONS:\n{agent_results['treatment']}\n\n"

        source_context = ("The following sources from Ayurvedic texts were "
                          "consulted by the agents:\n")
        for i, source in enumerate(retrieved_sources, 1):
            metadata = source.get('metadata', {})
            source_ref = (f"{metadata.get('section', 'Unknown Section')} - "
                          f"{metadata.get('chapter', 'Unknown Chapter')}")
            source_context += f"- Source {i}: {source_ref}\n"
\end{lstlisting}

The synthesis system prompt does two jobs. First, a silent review pass that is explicitly excluded from the output:

\begin{lstlisting}[style=plain]
**BACKGROUND VERIFICATION (DO NOT print in output):**
1. Clinical Consistency (Yukti): Ensure the Vikriti analysis aligns logically
   with the Chikitsa plan. A diagnosis of 'Ama' accumulation must be followed
   by 'Deepana-Pachana' recommendations. A Kapha-dominant Vikriti should not
   have Kapha-aggravating foods in the diet plan.
2. Textual Accuracy (Shastra): If a specific text was referenced, confirm the
   recommendations are consistent with that classical source.
3. Safety (Ahimsa): Double-check that all Shodhana procedures include a clear
   and non-negotiable warning requiring qualified supervision.
\end{lstlisting}

Second, a fixed four-section report skeleton:

\begin{table}[H]
\centering
\small
\begin{tabularx}{\textwidth}{@{}c l X@{}}
\toprule
\textbf{\#} & \textbf{Heading} & \textbf{Required content} \\
\midrule
1 & Provisional Diagnosis & Prakriti baseline, then Vikriti with dominant doshas, Gunas, Samprapti, affected Dhatus and Srotas \\
2 & Treatment Protocol & Chikitsa Sutra, Shodhana with safety warnings, Shamana split into classical formulations and single herbs \\
3 & Dietary \& Lifestyle & Pathya and Apathya lists, each with justification \\
4 & Clinical Summary & Practitioner-to-practitioner precis \\
\bottomrule
\end{tabularx}
\caption{The report structure the synthesis prompt enforces. Section numbering and headings are dictated verbatim.}
\end{table}

The call itself runs cooler and longer than any agent call:

\begin{lstlisting}[style=py,firstnumber=147]
        return self.llm_client.generate(prompt=synthesis_prompt,
                                        system_prompt=system_prompt,
                                        temperature=self.temperature,   # 0.2
                                        max_tokens=4000,
                                        conversation_history=conversation_history)
\end{lstlisting}

\section{A full slow-path trace}

Putting it together for the query \emph{``I have acid reflux, skin rashes, and a short temper. What should I do?''}:

\begin{figure}[H]
\centering
\begin{tikzpicture}[x=1mm,y=1mm,font=\scriptsize]
  % lifelines
  \foreach \x/\n in {0/UI, 28/Orch, 60/Agents, 96/Chroma, 128/LLM} {
    \node[blk=slate, minimum width=20mm, font=\scriptsize] at (\x,0) {\n};
    \draw[lifeline] (\x,-4) -- (\x,-104);
  }

  % messages
  \draw[arb=clay]      (0,-10)  -- (28,-10)  node[midway,above,font=\tiny]{query};
  \node[right,font=\tiny\itshape,color=slate] at (30,-16) {analyze\_query -> \{p:T, d:F, t:T\}};

  \draw[arb=deepgreen] (28,-24) -- (60,-24) node[midway,above,font=\tiny]{prakriti.process};
  \draw[arb=plum]      (60,-30) -- (96,-30) node[midway,above,font=\tiny]{where category=prakriti};
  \draw[arb=plum,dashed](96,-36) -- (60,-36) node[midway,above,font=\tiny]{5 chunks};
  \draw[arb=brick]     (60,-42) -- (128,-42) node[midway,above,font=\tiny]{prompt + context, T=0.3};
  \draw[arb=brick,dashed](128,-48) -- (60,-48) node[midway,above,font=\tiny]{Pitta-dominant};
  \draw[arb=deepgreen,dashed](60,-54) -- (28,-54);

  \draw[arb=deepgreen] (28,-62) -- (60,-62) node[midway,above,font=\tiny,align=center]{treatment.process\\+ Prakriti};
  \draw[arb=plum]      (60,-68) -- (96,-68) node[midway,above,font=\tiny]{where category=treatment};
  \draw[arb=plum,dashed](96,-74) -- (60,-74) node[midway,above,font=\tiny]{5 chunks};
  \draw[arb=brick]     (60,-80) -- (128,-80) node[midway,above,font=\tiny]{prompt + context, T=0.4};
  \draw[arb=brick,dashed](128,-86) -- (60,-86);
  \draw[arb=deepgreen,dashed](60,-90) -- (28,-90);

  \draw[arb=clay]      (28,-96) -- (128,-96) node[midway,above,font=\tiny]{synthesis, T=0.2, 4000 tok};
  \draw[arb=clay,dashed](128,-101) -- (0,-101) node[midway,above,font=\tiny]{four-section report};
\end{tikzpicture}
\caption{Slow path with two agents active. The dosha agent is skipped because no word in the query matches its trigger list --- ``acid reflux'' and ``rashes'' are not in the vocabulary, though a clinician would call them symptoms.}
\end{figure}

That trace also shows the routing weakness concretely. The query describes symptoms, but the word \emph{symptom} does not appear, so the diagnostic agent never runs and the treatment plan is built on a constitutional assessment with no imbalance analysis behind it.

\section{Cost of one consultation}

\begin{table}[H]
\centering
\small
\begin{tabular}{@{}l c c c@{}}
\toprule
\textbf{Path} & \textbf{Retrievals} & \textbf{LLM calls} & \textbf{Requested output tokens} \\
\midrule
Fast path & 0 & 1 & 2,000 \\
One agent + synthesis & 1 & 2 & 3,000 + 4,000 \\
Two agents + synthesis & 2 & 3 & 6,000 + 4,000 \\
All three + synthesis & 3 & 4 & 9,000 + 4,000 \\
\bottomrule
\end{tabular}
\caption{Per-query cost. Calls are sequential, not parallel --- the cascade requires it, since each agent's input includes the previous one's output.}
\end{table}

The sequential dependency is real for dosha and treatment, which consume upstream findings. Prakriti and dosha could in principle run concurrently against a query where the treatment agent is inactive, but the code does not attempt it.
